\documentclass[11pt,a4paper]{article}
\usepackage[margin=25mm,headheight=15pt]{geometry}
\usepackage[T1]{fontenc}
\usepackage[utf8]{inputenc}
\usepackage{lmodern,microtype}
\usepackage{amsmath,amssymb,amsthm}
\usepackage{booktabs,longtable,tabularx,array,enumitem}
\usepackage{xcolor,listings,xurl,hyperref,fancyhdr,needspace}
\definecolor{ink}{HTML}{163247}
\definecolor{accent}{HTML}{146B75}
\definecolor{pale}{HTML}{F2F6F8}
\hypersetup{colorlinks=true,linkcolor=accent,urlcolor=accent,citecolor=accent,
  pdftitle={Four Further Lean Proof Simplifications for DiophantineConsistency},
  pdfauthor={Research assistance for Vladimir Reshetnikov}}
\pagestyle{fancy}
\fancyhf{}
\fancyhead[L]{\small\color{ink}ProveIt: further Lean proof simplifications}
\fancyhead[R]{\small 26 October 2026}
\fancyfoot[C]{\small\thepage}
\renewcommand{\headrulewidth}{0.3pt}
\setlength{\parskip}{5pt}
\setlength{\parindent}{0pt}
\widowpenalty=10000
\clubpenalty=10000
\setlist{itemsep=3pt,topsep=4pt}
\lstset{basicstyle=\ttfamily\small,backgroundcolor=\color{pale},
  frame=single,rulecolor=\color{pale},breaklines=true,columns=fullflexible,
  keepspaces=true,showstringspaces=false,xleftmargin=3pt,xrightmargin=3pt}
\newcommand{\N}{\mathbb N}
\newcommand{\R}{\mathbb R}
\newcommand{\code}[1]{\nolinkurl{#1}}
\newcommand{\DC}{\code{Computability/DiophantineConsistency/Lean/DiophantineConsistency/}}
\newcommand{\Hten}{\code{Computability/HilbertTenthProblem/Lean/Diophantine/}}
\newcommand{\repo}{https://github.com/VladimirReshetnikov/ProveIt}
\newcommand{\rev}{9bd29079d5e1c6cc1c0bcaee6cb57d59b835cba9}
\newcommand{\src}[2]{\href{https://github.com/VladimirReshetnikov/ProveIt/blob/9bd29079d5e1c6cc1c0bcaee6cb57d59b835cba9/#1}{#2}}
\newtheorem{theorem}{Theorem}[section]
\newtheorem{lemma}[theorem]{Lemma}
\newtheorem{corollary}[theorem]{Corollary}
\newcolumntype{Y}{>{\raggedright\arraybackslash}X}
\newcolumntype{P}[1]{>{\raggedright\arraybackslash}p{#1}}
\begin{document}
\hypersetup{pageanchor=false}
\begin{titlepage}
\vspace*{12mm}
{\color{accent}\large RESEARCH AND IMPLEMENTATION NOTE}\par
\vspace{9mm}
{\color{ink}\Huge\bfseries Four further Lean\\[4pt]proof simplifications}\par
\vspace{6mm}
{\LARGE DiophantineConsistency after the earlier proposals}\par
\vspace{10mm}
{\large Prepared for Vladimir Reshetnikov\par 26 October 2026}\par
\vspace{13mm}
\begin{tabularx}{\textwidth}{@{}lY@{}}
\toprule
Reviewed repository & \href{https://github.com/VladimirReshetnikov/ProveIt}{VladimirReshetnikov/ProveIt}\\
Pinned revision & \code{9bd29079d5e1c6cc1c0bcaee6cb57d59b835cba9}\\
Lean version & \code{leanprover/lean4:v4.32.0}\\
Exclusion corpus & Six placed packets and five incoming ZIP archives\\
Deliverable & Mathematical derivations, standalone checked Lean prototypes, source and novelty evidence, reproduction scripts\\
\bottomrule
\end{tabularx}\par
\vspace{8mm}
\textbf{Principal recommendation.} Replace the exponential comparison in
\code{main_index_gt_first} by monotonicity of the positive Pell pair recurrence.
For a larger simplification, combine a natural-number perturbation estimate with
exact binomial quotient recovery. These changes remove specific mathematical
detours that remain after the earlier refactors.

\textbf{Verification boundary.} The accompanying small Lean modules are checked
with the repository's pinned compiler. They establish the reusable mathematical
and compilation principles. Production adapters, a full repository build,
exact supplied-program preservation, net line savings, and build-time changes
remain integration work. None of those unperformed steps is presented as a result.
\vfill
{\small Novelty means ``not found in the directly audited six packets and five
current incoming archives.'' It is not a claim of priority over all Lean or
mathematical literature. All earlier ideas are treated as already applied.}
\end{titlepage}
\hypersetup{pageanchor=true}
\pagenumbering{roman}
\tableofcontents
\clearpage
\pagenumbering{arabic}

\section{Findings and scope}

Four directions survive a fairly aggressive exclusion test. Three concern the
mathematics of the Pell argument; the fourth removes duplicated rule-schedule
definitions. They are deliberately unequal in scope: a small proof replacement,
a shared arithmetic kernel, direct reuse of an existing integer floor theorem, and
a maintenance-oriented compiler abstraction.

\begin{tabularx}{\textwidth}{@{}P{7mm}P{40mm}YP{31mm}@{}}
\toprule
 & Proposal & What disappears & Assessment\\
\midrule
1 & Compare Pell bases directly & Four exponential comparisons, predecessor exponents, and small-index case work in one index theorem & Best first pilot\\
2 & A natural perturbation and error-budget theorem & Repeated real division, reciprocal, positivity, and tolerance calculations in two upper estimates & Strongest consolidation\\
3 & Use the existing exact binomial quotient & A real fractional-part rounding detour and repeated cross-domain adapters & Strong companion to 2\\
4 & One rule-schedule syntax & Separate numeric and supplied-PR copies of each branch expression & Maintenance investment\\
\bottomrule
\end{tabularx}

The maintained directory contains 402 immediate Lean modules at this revision.
This count excludes its separate \code{checks} directory and the top-level
umbrella import. It describes the scale of the inspected tree; it does not imply
that every theorem was individually audited. The concrete targets below were
read in detail, and the evidence scripts record their exact bytes.

The user asks for simplifications that remain after previous proposals have
been applied. Accordingly, current source lengths are used only to identify
and size targets. They are not a counterfactual measurement of a repository
with all previous packets integrated. A new helper's own source and adapter
costs must be counted before a claim of net compression is made.

\subsection{The pinned baseline}

The inspected \code{main} head is \code{9bd29079d5e1}, whose commit metadata
records 26 October 2026. The root \code{lean-toolchain} pins Lean 4.32.0.
The root lockfile pins Mathlib to
\code{81a5d257c8e410db227a6665ed08f64fea08e997}.
The relevant public recurrence API was checked in that Mathlib revision,
rather than inferred from a different release. See references
\ref{ref:baseline} and \ref{ref:mathlib}.

All repository links in the report use the pinned revision. No changes were
pushed to the repository. The archive contains proposed reusable additions and
evidence, rather than a claim that the maintained development was refactored.

\section{Excluding earlier ideas}

The direct exclusion corpus consists of packets \code{2026-10-26a} through
\code{2026-10-26f} under \code{research/lean-proof-simplification}, plus all
five ZIPs currently in \code{docs/incoming}. The audit manifest records 149
placed text sources and all 400 members of the five ZIPs. Reports, source notes,
novelty audits, proposed interfaces, and relevant Lean artifacts were surveyed.
The archive contents were handled as inert data; none of their supplied scripts
was executed.

The audit includes ideas mentioned as optional or unselected directions. This
matters: an additive-equivalence treatment of aligned cells and an algebraic
interpretation of cyclic word packing were already discussed in evidence notes,
even though they were not prominent report headlines. The historical
\code{docs/incoming/README.md} was inspected for scope and routing; this report
does not claim to reread every unrelated manuscript transitively linked there.

\begin{tabularx}{\textwidth}{@{}P{42mm}Y@{}}
\toprule
Excluded family & Why it is not proposed again\\
\midrule
Decoder graph and fused depth/quotation proofs & Already present in several incoming packets, including the latest new-directions packet.\\
Signed congruence orbits and \code{natAbs} transport & Explicitly proposed in new-ideas and new-directions.\\
History folds, scanners, PR iteration, finite dispatch & Covered by the original packets, with supplied-program and cost concerns already discussed.\\
Circuit invariants, common tails, residual squares, depth certificates & Extensively covered, including the shared 99-node tail and generated certificates.\\
Witness layouts, source records, arithmetic row extraction & Existing proposals already isolate these transport and representation boundaries.\\
Coefficient envelopes, support bounds, finite pushforwards & Already covered as mathematical APIs and certificate schemes.\\
Canonical carries, explicit least/greatest cyclic carries & Already developed in additional-simplifications and new-directions.\\
Aligned-cell subgroup or additive equivalence & Explicitly identified in an earlier novelty audit as a refinement of old advice.\\
Internal constructor views, unique existence, graph fusion & Already present in the latest object-theory proposals.\\
\bottomrule
\end{tabularx}

The detailed \code{evidence/prior_art_audit.md} contains 56 substantive
exclusion families, smaller local suggestions, exact archive/member pointers,
and decisions on rejected candidates. \code{evidence/novelty_matrix.csv}
records the nearest earlier proposal for each retained direction.

\section{Compare Pell bases through the positive pair recurrence}
\label{sec:pell}

\subsection{The current detour}

The target is \code{NativePellBootstrap.main_index_gt_first}, lines 309--333
of the pinned file (reference \ref{ref:bootstrap}). It assumes
\[
 c=y_A(p),\qquad k=2y_P(n),\qquad k<c,
 \qquad 2A<2P-1,
\]
as well as positive-index and lower-bound hypotheses, and proves \(n<p\).
Here \(y_A\) denotes the second Pell coordinate with base \(A>1\).

The proof first excludes \(p=1\), introduces \(p-1\) and \(n-1\), and chains
\[
 c\le (2A)^{p-1}\le (2A)^{n-1}
 <(2P-1)^{n-1}\le y_P(n)\le k.
\]
This works, but compares approximations to the sequences when the sequences
themselves admit a simpler order comparison.

\subsection{The stronger mathematical interface}

Use the simultaneous positive recurrence
\begin{align*}
 x_A(0)&=1,&y_A(0)&=0,\\
 x_A(j+1)&=A x_A(j)+(A^2-1)y_A(j),&
 y_A(j+1)&=x_A(j)+A y_A(j).
\end{align*}

\begin{lemma}[Monotonicity in the base]
If \(1<A\le B\), then for every \(j\in\N\),
\[
 x_A(j)\le x_B(j),\qquad y_A(j)\le y_B(j).
\]
\end{lemma}
\begin{proof}
At zero the pairs coincide. For the successor step, all coordinates and
coefficients are nonnegative, \(A\le B\), and \(A^2-1\le B^2-1\).
Monotonicity of natural addition and multiplication transports both induction
hypotheses through the displayed recurrence.
\end{proof}

The simultaneous form is useful. The second-order recurrence
\(y_A(j+2)=2A y_A(j+1)-y_A(j)\) contains a subtraction whose monotonicity points
in the wrong direction for a direct comparison. The positive pair recurrence
avoids that issue completely.

\begin{theorem}[Index comparison from an actual value gap]
Suppose \(1<A\le P\), \(c=y_A(p)\), \(y_P(n)\le k\), and \(k<c\).
Then \(n<p\).
\end{theorem}
\begin{proof}
If \(p\le n\), base monotonicity and index monotonicity give
\[
 c=y_A(p)\le y_P(p)\le y_P(n)\le k<c,
\]
which is impossible.
\end{proof}

The stronger helper does not need \(p>0\), \(n>0\), or \(c>1\).
It uses \(A\le P\) in place of the old strict exponential-base separation.
The current public theorem can keep its exact signature: infer \(A\le P\)
from its old hypothesis, infer \(y_P(n)\le k\) from \(k=2y_P(n)\), and apply
the helper. This keeps downstream changes small while improving the actual proof.

\subsection{Lean boundary and payoff}

\code{lean/ArithmeticPellMonotonicity.lean} proves the pair recurrence comparison
using only \code{Std}. It also proves an adapter for any two coordinate
functions satisfying the same zero and successor equations, index monotonicity,
and the gap theorem. The base-monotonicity theorem has an empty axiom list
under \texttt{\#print axioms}; the recurrence adapter reports \code{propext}.

Pinned Mathlib supplies \code{Pell.xn_succ}, \code{Pell.yn_succ}, and
\code{Pell.strictMono_y}. Its coefficient \code{Pell.d} is private; an adapter
can expose \(A*A-1\) by definitional equality with \code{change}, without
referencing a private name. The exact ProveIt adapter is still to be compiled.

The immediate gain is a proof that follows the mathematical reason for the
index ordering and a weaker reusable interface. The original target is only
25 physical lines, so adding a general helper may not reduce the first client's
total line count. It is still the safest pilot: the transformation is local,
its core is checked, and the retained wrapper makes acceptance straightforward.
Deleting other base-separation lemmas would require a separate dependency check.

\section{A natural-number perturbation kernel for both Pell upper bounds}
\label{sec:perturbation}

\subsection{What remains duplicated}

\code{NativeHalfBinomial.upper_estimate_scaled} and
\code{Jones1980.Ratio90.upper_estimate'} contain nearly the same
growth comparison. The former allows the normalized ratio to be below
\(2(Y+1)\) and concludes an error below one; the latter uses \(Y+1\) and
concludes an error below one half. Their current theorem regions span
91 physical lines each. Those 182 lines are the target's
present size, not a measured deletion budget (references \ref{ref:half} and
\ref{ref:ratio}).

The two proofs introduce real quantities, prove positivity, clear fractions,
bound a perturbed power, cancel a positive power, and only then specialize the
error tolerance. Earlier advice to replace trivial existential bindings with
\code{set} improves presentation, but leaves this mathematical duplication.

\subsection{An exact replacement for the perturbed-power estimate}

\begin{theorem}[Natural perturbation]
For \(H,b,r\in\N\), if \(2rb\le H\), then
\begin{equation}
 H(H+b)^r\le H^r(H+2rb).
 \label{eq:perturb}
\end{equation}
\end{theorem}
\begin{proof}
Induct on \(r\), retaining its bound as a hypothesis. At zero both sides equal
\(H\). At a successor, apply the induction hypothesis and multiply by \(H+b\).
The remaining comparison is
\[
 (H+2rb)(H+b)\le H\bigl(H+2(r+1)b\bigr).
\]
After expanding and canceling common summands, this is exactly
\(b(2rb)\le bH\), supplied by \(2rb\le H\). The successor's stronger
assumption implies this bound and the induction hypothesis's assumption.
All terms are natural numbers; no division or subtraction is needed.
\end{proof}

For \(H>0\), equation \eqref{eq:perturb} is the denominator-cleared form of
\((1+b/H)^r\le 1+2rb/H\). Its direct induction avoids the real reciprocal
argument in the current \code{one_add_pow_le} proof.

\subsection{A reusable power-sandwich theorem}

\begin{theorem}[Transfer the perturbation through growth bounds]
Let \(H>0\). Suppose
\[
 ZH\le Q(H+b),\qquad V\le Z^r,\qquad Q^r\le W,
 \qquad 2rb\le H.
\]
Then
\begin{equation}
 VH\le W(H+2rb).
 \label{eq:sandwich}
\end{equation}
\end{theorem}
\begin{proof}
The exact chain is
\begin{align*}
 VH H^r
 &\le H(ZH)^r\\
 &\le H Q^r(H+b)^r\\
 &\le Q^r H^r(H+2rb)\\
 &\le W H^r(H+2rb).
\end{align*}
Cancel the positive factor \(H^r\). This uses \eqref{eq:perturb} once.
\end{proof}

This theorem is useful because it exposes precisely what each application must
provide. It does not package two existing proofs behind a common record.
There is one induction and one comparison argument.

\subsection{The exact Pell instantiation}

Use the following names, with \(U,Y,r\ge1\):
\begin{align*}
 H&=Y(U+1),& P&=2UY^2+1,\\
 D&=U^r,& T&=(U+1)^{2r},\\
 K&=y_P(r+1),& c&=y_{H+2}(2r+1),\\
 Z&=(2(H+2))^2U,& Q&=(2P-1)(U+1)^2.
\end{align*}
The existing natural \code{upper_key} theorem gives
\(ZH\le Q(H+9)\) when \(16\le H\).
The existing Pell growth bounds give
\[
 cD\le Z^r,\qquad Q^r\le KT.
\]
Apply \eqref{eq:sandwich} with \(b=9\), \(V=cD\), and \(W=KT\).
Under \(18r\le H\), the common conclusion is
\begin{equation}
 \boxed{cDH\le KT(H+18r).}
 \label{eq:majorant}
\end{equation}
Thus the shared estimate can remain entirely in \(\N\).
The existing stronger hypothesis \(72r<U+1\) readily supplies the perturbation
condition, but the core need not demand that stronger threshold.

For interpretation, dividing \eqref{eq:majorant} by positive denominators yields
\[
 \frac cK-\xi\le\frac{18r\xi}{Y(U+1)},
 \qquad \xi=\frac{T}{D}.
\]
This real inequality is a readable consequence of an integer certificate.
It need not be the implementation interface.

\subsection{Make the error tolerance a separate budget}

The following natural-number implication is checked in
\code{ArithmeticPerturbation.lean}:
\begin{equation}
 \left.
 \begin{aligned}
 cDH&\le KT(H+e),\\
 qTe&<DH,\qquad K>0
 \end{aligned}
 \right\}
 \quad\Longrightarrow\quad
 qcD<K(qT+D).
 \label{eq:budget}
\end{equation}
For positive \(q,D,K\), the conclusion says that the error in \(c/K\)
relative to \(T/D\) is below \(1/q\). Its proof multiplies the majorant by
\(q\), multiplies the budget by \(K\), adds the inequalities after distributing,
and cancels the common factor \(H\). The budget itself ensures the relevant
product is positive.

For the old half-error estimate choose \(q=2\) and \(e=18r\).
For the native estimate choose \(q=1\), then use its actual denominator
\(k=2K\). The checked native wrapper concludes
\begin{equation}
 2cD<k(T+D),
 \label{eq:native-upper}
\end{equation}
which is exactly the cross-multiplied upper inequality needed for the next
section's rounding theorem.

The familiar assumption \(72r<U+1\) gives
\[
 36r(Y+1)\le72rY<Y(U+1)=H.
\]
Consequently \(T<D(Y+1)\) supplies the budget \(36rT<DH\) for the old
half-error theorem, and \(T<2D(Y+1)\) supplies \(18rT<DH\) for the native
theorem. The latter bound on \(T\) already follows from the existing lower
estimate and the source interval by multiplying and canceling positive
natural numbers. No additional analytic estimate is required.

\subsection{Integration and expected effect}

The common kernel belongs below both consumers, beside \code{upper_key} or
in a small module imported by \code{PellRatio90}. The native file already
imports that module. Keep the current real-valued public theorems as
compatibility wrappers, proved by one denominator-clearing conversion if
clients still use their old statements.

The current \code{upper_estimate'} proof does not use its
\code{hYr : U ^ r <= Y} parameter after the declaration. The new core can omit
it; retaining it on the old wrapper preserves compatibility.

The checked artifacts establish the natural perturbation, the abstract
power-sandwich transfer, and the error-budget conversion. Connecting the
already available Pell growth bounds and \code{upper_key} to their exact
parameters is still a production adapter. This is the most promising
consolidation because the shared argument is visibly present in two current
proofs, and the new integer interface serves both.

\section{Recover the half-binomial value through an exact integer quotient}
\label{sec:rounding}

\subsection{The mathematical quantity to expose}

The current \code{half_binomial_of_ratio} obtains an integer part and a real
fractional remainder from \code{Jones1982.rho_split_c}, proves the integer
part even, compares two real intervals, then casts inequalities back to
naturals to identify the integer. The target lines are 146--177 of
\code{NativeHalfBinomial.lean}; \code{source_half_binomial}, at lines
193--242, builds its real arguments mainly at lines 210--241
(reference \ref{ref:half}).

Write \(T=(X+1)^{2r}\), \(D=X^r\), and \(C_r=\binom{2r}{r}\).
The needed natural-number certificate is already available in the imported
\code{Jones1980.Ratio90.binomial_floor_two}. Under the same native bound on \(X\),
it returns \(Q,L,W\) with
\[
 T=QD+L,\qquad 4L<D,\qquad Q=C_r+WX.
\]
The change is to use that existing interface directly and apply an integer
rounding theorem. It requires no new development of the binomial expansion.
The following derivation explains the existing certificate's mathematics;
it is not presented as a new binomial identity.

For clarity, consider the exact Euclidean decomposition of the numerator.
For \(r\ge1\), put
\[
 T=(X+1)^{2r},\qquad D=X^r,\qquad C_r=\binom{2r}{r}.
\]
Splitting the binomial sum at its middle term gives the explicit naturals
\begin{align*}
 L&=\sum_{j=0}^{r-1}\binom{2r}{j}X^j,\\
 W&=\sum_{j=0}^{r-1}\binom{2r}{r+1+j}X^j
\end{align*}
and the identity
\begin{equation}
 T=D(C_r+XW)+L.
 \label{eq:split}
\end{equation}

\Needspace{9\baselineskip}
\begin{lemma}[The lower binomial block is a remainder]
If \(X>2^{2r}\) and \(r\ge1\), then \(0<L<D\).
\end{lemma}
\begin{proof}
The term at \(j=0\) is one. For each \(j<r\),
\(X^j\le X^{r-1}\), and the partial sum of binomial coefficients is at most
the full sum \(2^{2r}\). Thus
\[
 L\le2^{2r}X^{r-1}<XX^{r-1}=X^r=D.
\]
The strict inequality uses positivity of \(X^{r-1}\), including when \(r=1\).
\end{proof}

The current stronger source bound \(2^{2r+1}\le X\) implies this condition.
For \(r\ge1\), the identity
\(\binom{2r}{r}=2\binom{2r-1}{r-1}\) makes \(C_r\) even. If \(X\) is even,
then \(C_r+XW=2t\) for a natural \(t\). Equation \eqref{eq:split} becomes
\begin{equation}
 \boxed{T=2tD+L,\qquad 0\le L<D.}
 \label{eq:half-quotient}
\end{equation}
Equivalently, \(T/D=2t\) in natural division, and its remainder is \(L\).
The existing \code{binomial_floor_two} supplies a stronger remainder bound
than needed here. Its \(4L<D\) immediately implies \(L<D\).
Alternatively, \code{Jones1982.lemma_2_24} with its scale parameter set to one
already gives the required natural decomposition. Neither existing theorem
needs to be reproved. The remaining adapter applies that theorem, establishes
the already-used parity fact, and connects it to the checked rounding lemma.

\subsection{A checked rounding theorem with integer inputs}

\begin{theorem}[Half-quotient rounding]
Let \(D,k>0\), \(0\le L<D\), and \(T=2tD+L\). Suppose
\begin{align}
 kT&<2cD<k(T+D),\label{eq:ratio-window}\\
 kY&<c<k(Y+1).\label{eq:source-window}
\end{align}
Then \(Y=t\), and consequently \(2Y=C_r+XW\) in the binomial application.
\end{theorem}
\begin{proof}
First show \(kt<c<k(t+1)\). If \(c\le kt\), then
\(2cD\le2ktD\le k(2tD+L)=kT\), contradicting the lower inequality.
If \(k(t+1)\le c\), then
\[
 2cD\ge2k(t+1)D>k(2tD+L+D)=k(T+D),
\]
contradicting the upper inequality. The strict step uses \(L<D\) and \(k>0\).
Finally, two strict intervals of the form \((kj,k(j+1))\) containing the same
integer \(c\) have the same index. Apply this to \(j=t\) and \(j=Y\).
\end{proof}

\code{ArithmeticQuotientRounding.lean} checks this theorem for arbitrary
natural inputs. It includes the interval uniqueness lemma and does not assume
a Pell interpretation or any real-number theorem. The remainder is allowed
to be zero, although the concrete binomial remainder is positive.

\subsection{How the existing source supplies the premises}

The current lower theorem gives \(kT<2cD\) directly over naturals.
Proposal 2 supplies the upper half of \eqref{eq:ratio-window} in precisely
the form \eqref{eq:native-upper}. The literal source equations give
\[
 c=kY+\eta,\qquad k=\eta+\zeta,
\]
with positive \(\eta,\zeta\), and therefore \eqref{eq:source-window}.
Once the existing binomial certificate and parity are supplied, the source theorem can apply
the checked rounding result and return the same witness \(W\).

The numerical identity being proved is unchanged. In particular, the factor
two in the native first Pell coordinate remains explicit throughout. As a
small exact illustration, take \(r=1\), \(X=8\). Then
\[
 T=81,\quad D=8,\quad C_r=2,\quad W=1,\quad L=1,\quad t=5.
\]
With \(k=16\), \(c=85\), and \(Y=5\), both windows hold and the recovered
identity is \(2Y=2+8\). This is an illustration of the local rounding lemma,
not a claim that these values form a complete source witness.

\subsection{The marginal contribution after the earlier packets}

Earlier Euclidean certificates package quotient/remainder constraints in
register and quadratic encodings. This proposal instead identifies the
\emph{particular integer quotient of a binomial expansion} and uses it to
replace the real rounding stage. The distinction is the eliminated proof
boundary, not the word ``Euclidean.''

Proposals 2 and 3 compose particularly well: the first ends in an integer
cross-product upper bound, and the second starts with exactly that bound.
They do not yet establish that the entire repository has no real-number
dependencies. Their target is the local Pell-to-binomial recovery chain.
The work is therefore smaller than introducing a fresh binomial-sum library.
The adapter and shared rounding helper still have a cost; source reduction
is an integration measurement, not a consequence of the ordinary proof alone.

\section{Define rule schedules once and derive both interpretations}
\label{sec:schedules}

\subsection{The duplication that finite dispatch does not remove}

In \code{AnnotatedRuleSchedules.lean}, the numerical definitions
\code{outputBody} and \code{premiseBodies} describe every rule branch.
The supplied-program definitions \code{P.outputBodyProgram} and
\code{P.premiseBodyProgram} repeat those expressions using realizers.
The four table regions contain 63 nonblank source lines. The existing
\code{ArithmeticFiniteSwitch.program} already handles dispatch; it cannot
remove duplication inside the branch bodies (reference \ref{ref:schedules}).

The four independent bridges to the original rule calculus occupy only
13 lines together. They are \code{output_quote}, \code{premises_quote},
\code{ofRule_output}, and \code{ofRule_premises}. They should remain compact
independent checks. Earlier advice already cautioned against replacing such
short finite proofs with metaprogramming.

\subsection{A small syntax is sufficient}

The actual schedules need ten named coordinate fields:
\code{context}, \code{shifted}, \code{a}, \code{b}, \code{c},
\code{instI}, \code{instJ}, \code{shiftC}, \code{i}, and \code{j}.
A term language needs only field access, natural literals, atom construction,
fork construction, and list cons. Binary/unary formula constructors, variable
nodes, and bottom are transparent abbreviations in this language.

Store one \code{Schedule} per tag, consisting of an output term and a list of
premise terms. For example, implication elimination has output
\(\operatorname{fork}(G,b)\) and premises
\[
 [\operatorname{fork}(G,\operatorname{binary}(3,a,b)),
   \operatorname{fork}(G,a)].
\]
The table has all 17 valid tags and the current zero-output/empty-premises
default for every invalid tag.

Numerical evaluation interprets the syntax in the existing coding operations.
A backend compiles it using the actual supplied PR constructors. Its laws
state that each primitive program constructor evaluates to the corresponding
numerical operation. Structural induction then proves
\begin{equation}
 \operatorname{run}(\operatorname{compile}(t),z)
 =\operatorname{eval}(t,\operatorname{coordinates}(z))
 \label{eq:compiler}
\end{equation}
for every raw input \(z\). This theorem does not require a valid coordinate
packet. Validity is used later, in the independent bridge to the original
rule's output and premises.

\subsection{What has actually been checked}

\code{StructuralSchedule.lean} includes the full 17-rule table, its invalid
default, independent copies of the legacy numeric tables, a generic backend,
and proofs of:
\begin{itemize}
\item \code{schedule_matches_legacy}: exact numeric table agreement for every
natural-number tag, including the invalid default;
\item \code{Backend.compile_correct}: equation \eqref{eq:compiler} for every
term and raw input;
\item \code{schedule_correct}: simultaneous correctness of the output and
the encoded premise list.
\end{itemize}
All three have empty axiom lists. The companion Python transcription checker
compares the independent legacy tables against the pinned production source
after a documented renaming of field access and coding operations. Both
comparisons pass. Thus the validation covers the actual table, not just a
generic language unrelated to its 17 cases.

\subsection{The important remaining data-preservation obligation}

Instantiation with \code{ScalarPRCode} must use the current projection
realizers and coding combinators. Before replacing maintained code, prove
equality of each compiled branch's executable \code{.val} with the original
branch tree. Equality of evaluated functions alone does not establish this.

The pinned source currently states node counts 2,537 and 2,377 for the output
program before and after optimization, and 2,683 and 2,483 for the premise
program. These numbers were read from source; the standalone prototype does
not rerun those theorems. Replaying them is a production acceptance gate.

Generic syntax and backend definitions have a cost. This first client could
grow in total source size even though the number of independent copies of
each branch expression falls from two to one. The existing dependent realizer
type already catches many semantic inconsistencies, so the benefit is chiefly
maintenance and a common correctness argument. Adopt this after the small Pell
pilot, and judge it by actual net size, change effort, and program-data
preservation rather than by an assumed speedup.

\section{Integration order and measurements}

The following order keeps each proposed change independently reviewable.

\begin{enumerate}
\item Add the Pell recurrence adapter and base-monotonicity theorem. Replace
the body of \code{main_index_gt_first} while retaining its header. Compile
that file and its actual clients.
\item Add the natural perturbation and power-sandwich helpers below both
upper-bound consumers. Derive the common majorant from the existing growth
lemmas and \code{upper_key}. Recover the two existing upper-bound statements
as wrappers before modifying clients.
\item Reuse \code{binomial_floor_two} and the existing parity fact. Combine the
majorant budget, existing lower estimate, literal source interval, and checked
rounding theorem in \code{source_half_binomial}.
\item Instantiate the schedule backend. Retain the independent bridges,
prove executable-tree equality, and replay the four count theorems.
\end{enumerate}

Measure production code plus new shared helpers plus adapters together.
Record the number of independent inductions and cross-domain transitions as
well as physical/nonblank lines. For the schedule change, record the number
of independent branch descriptions and whether program trees are unchanged.
Compile time and kernel time should be measured only after the exact adapters
are integrated. This packet reports no benchmark or project-wide reduction.

These are implementation acceptance criteria, rather than additional new
ideas. Earlier packets already discuss profiling, axiom audits, narrow
simplification, and regression checks; those are assumed available.

\section{Artifact guide and validation boundary}

The archive is self-contained apart from an installed Lean 4.32.0 compiler
and ordinary LaTeX/Python tools. It does not download a repository or install
dependencies when its check scripts are run.

\begin{tabularx}{\textwidth}{@{}P{62mm}Y@{}}
\toprule
Artifact & Purpose\\
\midrule
\code{report.tex}, \code{report.pdf} & Complete report and compiled rendering.\\
\code{lean/ArithmeticPellMonotonicity.lean} & Positive pair recurrence, base and index comparisons, recurrence adapter.\\
\code{lean/ArithmeticPerturbation.lean} & Natural perturbed-power inequality, abstract majorant transfer, exact tolerance selection, and native doubled-denominator wrapper.\\
\code{lean/ArithmeticQuotientRounding.lean} & Exact rounding from two integer windows.\\
\code{lean/StructuralSchedule.lean} & Actual 17 schedules, legacy-table agreement, generic compilation correctness.\\
\code{evidence/} & Full exclusion audit, candidate review, hashes, source targets, novelty matrix, and check outputs.\\
\code{scripts/} & Reproduction, source transcription, and exact arithmetic experiments.\\
\bottomrule
\end{tabularx}

All four standalone modules are checked with Lean 4.32.0; their compiler
outputs are included. No \code{sorry}, \code{admit}, added axiom declaration,
\code{native_decide}, or unsafe declaration is used. Some proofs report the
ordinary Lean foundational axioms \code{propext}, \code{Quot.sound}, and/or
\code{Classical.choice}; an empty axiom list is claimed only where the actual
audit output says so. ``No added axioms'' does not mean every theorem has an
empty dependency list.

The exact arithmetic experiments use integer calculations, including the
binomial decomposition and representative Pell majorants. They provide
additional checks on factors and signs; they do not substitute for the
unbounded proofs. The source transcription check establishes correspondence
with this pinned repository snapshot, not arbitrary future revisions.

The remaining formalization tasks are explicit: the actual Mathlib Pell
adapter, the source instantiation of the natural majorant, the adapter from
the existing binomial certificate to rounding, and the actual PR backend/data bridge. A full ProveIt
build was not performed. The report's source-size observations therefore
remain target measurements and its expected benefits remain proposals.

\section{References and exact source targets}

\begin{enumerate}[label={[\arabic*]},leftmargin=9mm]
\item\label{ref:baseline}
\src{lean-toolchain}{ProveIt toolchain},
\src{lake-manifest.json}{dependency lockfile}, and
\href{https://github.com/VladimirReshetnikov/ProveIt/commit/9bd29079d5e1c6cc1c0bcaee6cb57d59b835cba9}{reviewed commit}.

\item\label{ref:bootstrap}
\src{Computability/DiophantineConsistency/Lean/DiophantineConsistency/NativePellBootstrap.lean}{\code{NativePellBootstrap.lean}}:
\code{main_index_gt_first}, lines 309--333; associated classification and
source-base lemmas were inspected for context.

\item\label{ref:half}
\src{Computability/DiophantineConsistency/Lean/DiophantineConsistency/NativeHalfBinomial.lean}{\code{NativeHalfBinomial.lean}}:
\code{upper_estimate_scaled}, 18--108;
\code{doubled_upper_estimate}, 111--142;
\code{half_binomial_of_ratio}, 146--177;
\code{source_half_binomial}, 193--242.

\item\label{ref:ratio}
\src{Computability/HilbertTenthProblem/Lean/Diophantine/Paper1980/PellRatio90.lean}{\code{Paper1980/PellRatio90.lean}}:
\code{upper_key}, 189--195;
\code{one_add_pow_le}, 198--219;
\code{upper_estimate'}, 222--312; and the already available
\code{binomial_floor_two}, beginning at line 367.

\item\label{ref:lower}
\src{Computability/DiophantineConsistency/Lean/DiophantineConsistency/NativePellPower.lean}{\code{NativePellPower.lean}}:
\code{doubled_ratio_lower} and power recovery.
\src{Computability/HilbertTenthProblem/Lean/Diophantine/Common/Pell.lean}{\code{Common/Pell.lean}}:
\code{psi_succ_le_pow} and \code{pow_le_psi_succ} (the source spells psi with
the Greek symbol), the existing natural growth estimates.

\item\label{ref:binomial}
\src{Computability/HilbertTenthProblem/Lean/Diophantine/Paper1982/Ratio.lean}{\code{Paper1982/Ratio.lean}}:
\code{lemma_2_24}, beginning at line 44, the existing natural certificate;
and \code{rho_split_c}, its real-remainder interface used by the current
native rounding proof.

\item\label{ref:schedules}
\src{Computability/DiophantineConsistency/Lean/DiophantineConsistency/AnnotatedRuleSchedules.lean}{\code{AnnotatedRuleSchedules.lean}}:
numeric and supplied-program tables, the four independent bridges, and
node-count declarations at lines 264--267.

\item\label{ref:prior}
\src{Computability/DiophantineConsistency/research/lean-proof-simplification/README.md}{Placed-packet index} and
\src{Computability/DiophantineConsistency/research/lean-proof-simplification/ADOPTION.md}{adoption ledger}.
The ledger's adopted/pending/deferred distinctions do not change this report's
assumption that all earlier proposals have already been applied.

\item\label{ref:incoming}
\href{https://github.com/VladimirReshetnikov/ProveIt/tree/9bd29079d5e1c6cc1c0bcaee6cb57d59b835cba9/docs/incoming}{Current incoming archives}:
\code{diophantine-lean-simplification-2026-10-26(1).zip},
\code{lean-proof-simplification-review-2026-10-26.zip},
\code{diophantine-lean-new-ideas-2026-10-26.zip},
\code{diophantine-lean-additional-simplifications-2026-10-26.zip}, and
\code{diophantine-lean-new-directions-2026-10-26.zip}.
Exact member references and hashes appear in the evidence files.

\item\label{ref:mathlib}
\href{https://github.com/leanprover-community/mathlib4/blob/81a5d257c8e410db227a6665ed08f64fea08e997/Mathlib/NumberTheory/PellMatiyasevic.lean}{Pinned Mathlib \code{PellMatiyasevic.lean}}:
positive pair definition at 104--106, \code{xn_succ} at 133,
\code{yn_succ} at 137, and \code{strictMono_y} at 375.
The file was inspected for the needed API; no exhaustive Mathlib-wide
novelty claim is made.
\end{enumerate}

\end{document}
